\documentclass[crop=false]{standalone}
% --- ПРЕАМБУЛА ДЛЯ ПОДДЕРЖКИ РУССКОГО ЯЗЫКА И МАТЕМАТИКИ ---
\usepackage[T2A]{fontenc}
\usepackage[utf8]{inputenc}
\usepackage[russian]{babel}
\usepackage{amsmath}
\usepackage{geometry}
\usepackage{amssymb}
\usepackage{graphicx}
\usepackage{float}
\usepackage{import}
\geometry{a4paper, top=2cm, bottom=2cm, left=2cm, right=2cm}
\newcommand{\sidecomment}[1]{\marginpar{\raggedright\scriptsize\itshape #1}}
\newcommand{\iu}{\mathrm{i}}

\newcommand{\Def}{%
    \text{Def}%
    \hspace{-0.85em}% Сдвиг назад ровно на середину (подберите под шрифт)
    \raisebox{0.1ex}{/}% Черта
    \hspace{1em}% Возврат к концу слова + небольшой отступ
}

\renewcommand{\Re}{\operatorname{Re}}
\renewcommand{\Im}{\operatorname{Im}}

\ifstandalone
    \newcommand{\lecturebreak}{} % В отдельном файле лекции разрыв не нужен
\else
    \newcommand{\lecturebreak}{\clearpage} % В полной книге - начинаем с новой страницы
\fi

\newcounter{lecturenumber} % Создаем счетчик для нумерации лекций
\newcommand{\lecture}[2]{
    \lecturebreak
    \refstepcounter{lecturenumber} % Увеличиваем счетчик на 1
    \par\vspace{2em}\noindent % Отступ сверху и убираем абзацный отступ
    {\Large\bfseries % Делаем заголовок крупным и жирным
    Лекция \thelecturenumber: #1 % Выводим "Лекция N: Тема"
    \hfill % Растягиваем пробел, чтобы дата была справа
    \large\textit{#2}} % Выводим дату курсивом
    \par\vspace{1em} % Небольшой отступ снизу
    \addcontentsline{toc}{section}{Лекция \thelecturenumber: #1} % Добавляем в оглавление
    \par
}

\newcommand{\TwoCols}[2]{%
  \par\noindent % Начать с новой строки без отступа
  \begin{minipage}[t]{0.48\textwidth}
    #1
  \end{minipage}% <-- Важный '%' для удаления пробела
  \hfill % Растягивающийся пробел между колонками
  \begin{minipage}[t]{0.48\textwidth}
    #2
  \end{minipage}% <-- Важный '%' для удаления пробела
  \par%
}

\pagestyle{plain} % Включаем нумерацию страниц\
\begin{document}
\lecture{}{17.04}
\begin{theorem}[Гюйгенса-Штейнера]

\end{theorem}
\[J = J_C + M \begin{pmatrix}
    y_C^2 + z_C^2  & x_C y_C & x_C z_C \\ x_C y_C & x_C^2 + z_C^2 & x_C z_C \\ x_C z_C & y_C z_C & x_C^2 + y_C^2
\end{pmatrix}\]
\begin{proof}
    %fixme сделать зарисовку, фото в телефоне
    \[
    \begin{aligned}
        J_{C_x} &= \sum_{k=1}^{N} m \left((z_k - z_C)^2 + (y_k - y_C)^2\right) = \sum_{k=1}^{N} m_k (z_k^2 - 2 z_k z_C + z_C^2 + y_k^2 - 2 y_k y_C + y_C^2) \\
        &= \sum_{k=1}^{N} m_k (z_k^2 + y_k^2) - 2\sum_{k=1}^{N} m_k(z_k z_C + y_k y_C) \sum_{k=1}^{N} m_k (z_C^2 + y_C^2) \\
        &= J_X + M (z_C^2 + y_C^2) -2 M (z_C^2 + y_C^2) = J_x - M(z_C^2 + y_C^2)
    \end{aligned}
    \]
    \[
    \begin{aligned}
    J_{C_{x y}} &= \sum_{k=1}^{N} m_k (x_k - x_C) (y_k - y_C) = \sum_{k=1}^{N} m_k x_k y_k + \sum_{k=1}^{N} m_k x_C y_C - \sum_{k=1}^{N} m_k x_C y_k - \sum_{k=1}^{N} m_k x_k y_C \\
    &= J_x + x_C y_C M - M x_C y_C - M y_C x_C = J_x - M x_C y_C
    \end{aligned}
    \]
\end{proof}
Следствие: $\bar u_C$ — ось проходящая через центр масса, сонаправленная с $\bar u$
\[J_{\bar u} = J_{\bar u_C} + M d^2 \qquad d \text{ — расстояние между осями}\]
%fixme картинка с двумя осями, и расстоянием между ними

\noindent\keypoint \[J_{\bar u} = \bar u ^T J \bar u = \alpha^2 J_x + \beta^2 J_y + \gamma^2 J_z -2 \alpha \beta J_{xy} - 2 \beta \gamma J_{yc} - 2\alpha \gamma J_{xz}\]
\keypoint Положительно определённая квадратичная форма % fixme почему?

\Def Тогда существует ск, в которой она задаёт эллипсоид инерции
\[A x^2 + B y^2 + C z^2 = 1 \]
\Def Главная система координат — система координат в которой тензор инерции твёрдого тела имеет вид 
\[J = \begin{pmatrix}
    A & 0& 0\\0 & B & 0\\ 0&0& C 
\end{pmatrix}\]
Её оси будем называть главными осями твёрдого тела

\noindent \keypoint
\[\bar r^T J \bar r = 1 \qquad r^2 \bar u^T J \bar u = 1 \qquad r = \sqrt{\frac{1}{J_{\bar u}}} = \frac{1}{\sqrt{J_{\bar u}}}\]
То есть эллипсоид инерции будет более вытянутым в сторону оси с наименьшим моментом инерции

\section{Основные динамические величины в Кёниговой ск}
$O$ — неподвижная точка тела(в Кёниговой это центр масс)
\noindent \keypoint
Кинетический момент тела с неподвижной точкой
\[\bar K = J \bar \omega\]
\begin{proof}
\[
\begin{aligned}
\bar K &= \sum_{k=1}^{N} \bar r_k \times \bar v_k m_k = \sum_{k=1}^{N} \bar r_k \times (\bar \omega \times \bar r_k) m_k = \sum_{k=1}^{N} (\bar \omega r_k^2 - \bar r_k (\bar \omega, \bar r_k)) m_k \\
&= \sum_{k=1}^{N} (\bar \omega (x_k^2 + y_k^2 + z_k^2) - \bar r_k (\omega_x x_k + \omega_y y_k + \omega_z z_k)) m_k  \\
&= \begin{pmatrix}
    \sum_{k=1}^{N} \left( \omega _x ( x_k^2 + y_k^2 + z_k^2) - x_k  (\omega_x x_k + \omega_y y_k + \omega_z z_k) \right) m_K \\
    \sum_{k=1}^{N} \left( \omega _y ( x_k^2 + y_k^2 + z_k^2) - y_k  (\omega_x x_k + \omega_y y_k + \omega_z z_k) \right) m_K \\
    \sum_{k=1}^{N} \left( \omega _z ( x_k^2 + y_k^2 + z_k^2) - z_k  (\omega_x x_k + \omega_y y_k + \omega_z z_k) \right) m_K \\
\end{pmatrix} = \begin{pmatrix}
    \omega_x J_x - \omega_y J_{x y} - \omega_z J_{x z} \\
    -\omega_ x J_{x y} + \omega_y J_y - \omega_z J_{y z} \\
    -\omega_x J_{x z} - \omega_y J_{y z} + \omega_z J_z
\end{pmatrix} = J \bar \omega
\end{aligned}
\]
\end{proof}
\noindent \keypoint
Кинетическая энергия тела
\[T = \frac{1}{2} (\bar \omega , \bar K_O) = \frac{1}{2} J \omega^2\]
\begin{proof}
    \[T = \frac{1}{2} \sum_{k=1}^{N} m_k v_k^2 = \frac{1}{2} \sum_{k=1}^{N } m_k (\bar \omega \times \bar r_k)^2 = \frac{\omega^2}{2}  \underbracket{ \sum_{k=1}^{N} m_K(\bar u \times \bar r_k)^2}_{J_{\bar u}}  =  \frac{\omega^2 J_{\bar u}}{2} = \frac{1}{2}\bar \omega^T J \bar \omega = \frac{1}{2} \bar \omega^T \bar K_O = \frac{1}{2} (\bar \omega , \bar K_O)\]
\end{proof}

\section{Потенциальная система сил}
\Def Потенциальная система сил $\bar F_k$ — система сил, для которой существует скалярная функция $U$, называемая потенциальной
\[\bar F_k = \frac{\partial U}{\partial \bar r_k}  \qquad \Pi = -U \text{ — потенциальная энергия} \qquad U = U(\bar r_1 , \dots , \bar r_N)\]
\[\bar F_k = - \frac{\partial \Pi}{\partial \bar r_k}\]

\keypoint
Обобщённые силы потенциальной системы сил потенциальны
\[Q_i = \sum_{k=1}^{N} \bar F_k \frac{\partial \bar r_k}{\partial q_i} = \sum_{k=1}^{N} \frac{\partial U}{\partial \bar r_k} \frac{\partial \bar r_k}{ \partial q_i} = \frac{\partial U}{\partial q_i}\]

\keypoint
Дифференциал потенциальной энергии
\[d \Pi = \sum_{k=1}^{N} \frac{\partial \Pi}{\partial \bar r_k} d \bar r_k + \frac{\partial \Pi}{\partial t} dt\]

\keypoint \[d' A = \sum_{k=1}^{N} \bar F_k d \bar r_k = \sum_{k=1}^{N} \frac{\partial U}{\partial \bar r_k} d \bar r_k = - \sum_{k=1}^{N} \frac{\partial \Pi}{\partial \bar r_k} d \bar r_k = - d \Pi + \frac{\partial \Pi}{\partial t} dt\]

\keypoint для стационарных систем потеницальных сил
\[d' A = - d \Pi\]
\Def Консервативная система сил — мех с в которй действуют только стационарные потенциальные силы
\[d T = d' A = - d \Pi \implies d(T  +\Pi) = 0 \implies T + \Pi = \text{const}\]
\Def Полная энергия $E = T + \Pi$
\end{document}